\chapter{Matrix Algebra}
\label{ch:matrix}

\section{Why matrices?}

Recall that an $n\times 1$ vector $Y$ can collect $n$ observations of an
outcome, while an $n\times k$ matrix $X$ collects $k$ explanatory variables
for those same observations.  Most estimators in this course are formulas
built from $Y$ and $X$, so a compact set of matrix-algebra tools will save us
a great deal of work.

This chapter contains the tools needed in the main course.  Additional
algebra, including eigenvalues and matrix derivatives, is collected in
Appendix~\ref{app:matrix}.

\section{Scalars, vectors, and matrices}

A \emph{scalar} is a number.  A \emph{vector} is an ordered list of numbers,
usually written as a column.  A \emph{matrix} is a rectangular array of
numbers.  If $A$ has $n$ rows and $k$ columns, then it is an $n\times k$
matrix.  Its entry in row $i$ and column $j$ is $a_{ij}$ (or $A_{ij}$).
For example,
\[
  A=\begin{pmatrix}
    a_{11}&a_{12}&a_{13}\\
    a_{21}&a_{22}&a_{23}
  \end{pmatrix}
  =(A_1,A_2,A_3)
\]
is $2\times3$, where $A_j$ is its $j$th column.  A scalar is a $1\times1$
matrix, and a column vector is a matrix with one column.

\begin{definition}[Transpose]
The transpose $A'$ is obtained by interchanging rows and columns:
$(A')_{ij}=A_{ji}$.  If $A$ is $n\times k$, then $A'$ is $k\times n$.
\end{definition}

Thus, for the matrix above,
\[
  A'=\begin{pmatrix}
    a_{11}&a_{21}\\
    a_{12}&a_{22}\\
    a_{13}&a_{23}
  \end{pmatrix}.
\]

\begin{definition}[Square, symmetric, diagonal, and identity matrices]
A matrix is \emph{square} if it has the same number of rows and columns.  A
square matrix is \emph{symmetric} if $A=A'$.  A matrix is \emph{diagonal} if
all off-diagonal entries are zero.  The $k\times k$ identity matrix $I_k$ is
diagonal with every diagonal entry equal to one.
\end{definition}

Symmetric matrices appear repeatedly as covariance matrices and quadratic
forms.  Identity matrices play the same role for matrix multiplication that
the scalar $1$ plays for ordinary multiplication.

\section{Matrix multiplication}

Two matrices $A$ and $B$ are \emph{conformable} for the product $AB$ when the
number of columns of $A$ equals the number of rows of $B$.  If $A$ is
$m\times k$ and $B$ is $k\times r$, then $AB$ is $m\times r$, with
\begin{equation}
\label{eq:matrix-product}
  (AB)_{ij}=\sum_{\ell=1}^{k}a_{i\ell}b_{\ell j}.
\end{equation}
Always check dimensions before multiplying.

\subsection{Inner and outer products}

For two $k\times1$ vectors $a$ and $b$, the \emph{inner product} is the
scalar
\[
  a'b=\sum_{j=1}^{k}a_jb_j.
\]
The \emph{outer product} $ab'$ is instead a $k\times k$ matrix whose
$(i,j)$ entry is $a_i b_j$.  This distinction matters throughout
econometrics.  For instance, $X_iX_i'$ is a $k\times k$ outer product, while
$X_i'\beta$ is a scalar inner product.

If $X$ is $n\times k$ and its $i$th row is $X_i'$, then
\begin{equation}
\label{eq:xtx-sum}
  X'X=\sum_{i=1}^{n}X_iX_i'.
\end{equation}
This identity turns a sample average of outer products into compact matrix
notation.

\subsection{Rules worth remembering}

Matrix multiplication is generally not commutative: even when both products
exist, $AB$ need not equal $BA$.  It is associative and distributive whenever
the dimensions conform:
\[
  A(BC)=(AB)C,\qquad A(B+C)=AB+AC.
\]
Transpose reverses the order of a product:
\[
  (AB)'=B'A'.
\]

\begin{rcode}
A <- matrix(c(1, 2, 3, 4, 5, 6), nrow = 3, ncol = 2)
B <- matrix(c(2, 3, 4, 5, 6, 7), nrow = 3, ncol = 2)
t(A) %*% B

# The same (1,2) entry as an inner product of two columns
sum(A[, 1] * B[, 2])
\end{rcode}

\section{Linear independence, rank, and invertibility}

Vectors $a_1,\ldots,a_r$ are \emph{linearly independent} if
\[
  c_1a_1+\cdots+c_ra_r=0
\]
implies $c_1=\cdots=c_r=0$.  Otherwise at least one direction is redundant.

\begin{definition}[Rank and full rank]
The rank of a matrix is the dimension of its column space (equivalently, its
row space).  An $m\times k$ matrix has
\[
  \operatorname{rank}(A)\leq \min\{m,k\}.
\]
It is \emph{full column rank} if its rank is $k$, \emph{full row rank} if its
rank is $m$, and simply \emph{full rank} if its rank is
$\min\{m,k\}$.
\end{definition}

An inverse is defined only for a square matrix.  A $k\times k$ matrix $A$ is
invertible if there exists $A^{-1}$ satisfying
\[
  A^{-1}A=AA^{-1}=I_k.
\]
A square matrix is invertible if and only if it has rank $k$.  A rectangular
full-rank matrix is \emph{not} invertible, although one of its Gram matrices
may be:
\begin{itemize}
  \item $A'A$ is invertible if and only if $A$ has full column rank;
  \item $AA'$ is invertible if and only if $A$ has full row rank.
\end{itemize}
Both Gram matrices can be invertible simultaneously only when $A$ is square
and full rank.  Also,
\[
  \operatorname{rank}(A'A)=\operatorname{rank}(AA')
  =\operatorname{rank}(A).
\]
If conformable square matrices $A$ and $B$ are invertible, then
$(AB)^{-1}=B^{-1}A^{-1}$.

\section{Positive semidefinite matrices}

\begin{definition}[Positive definiteness]
A real symmetric $k\times k$ matrix $A$ is \emph{positive semidefinite}
(p.s.d.) if $x'Ax\geq0$ for every $x\in\mathbb R^k$.  It is \emph{positive
definite} (p.d.) if $x'Ax>0$ for every nonzero $x$.
\end{definition}

For any real $m\times k$ matrix $A$,
\[
  x'A'Ax=(Ax)'(Ax)=\lVert Ax\rVert^2\geq0,
\]
so $A'A$ is p.s.d.  It is p.d. precisely when $A$ has full column rank.
Similarly, $AA'$ is p.s.d. and is p.d. precisely when $A$ has full row rank.
This argument is both shorter and more general than appealing to an
eigendecomposition.

Covariance matrices are p.s.d. for the same reason.  If $W$ is a random
vector with finite second moments and covariance matrix $\Sigma$, then for
any $a$,
\[
  a'\Sigma a
  =\mathbb E\!\left[\{a'(W-\mathbb EW)\}^2\right]\geq0.
\]

The optional algebra in Appendix~\ref{app:matrix} gives detailed formulas,
eigenvalue characterizations for symmetric matrices, and the matrix
derivatives used to derive least squares.
